\chapter{The rule systems, and logarithmic differentiation}

\begin{orient}
\textbf{What this chapter is for.} When the algebra of the previous chapter runs out, you differentiate, and differentiating means choosing a rule. The choice is not free. The product, quotient and chain rules are adapted to three different shapes, and applying one to a shape it was not built for produces algebra that never ends: a correct answer buried under a cancellation you now have to perform by hand. This chapter gives explicit selection guidance, including the cases where the quotient rule is the wrong tool and should be replaced by a termwise split or by a negative power, and the generalised product rule for $n$ factors, which is what people reach for two-at-a-time and lose terms doing. It then develops logarithmic differentiation as the general solvent: take $\ln$ of both sides, differentiate implicitly, multiply back by $y$. That single manoeuvre handles many-factor products, the variable-base-variable-exponent form $f(x)^{g(x)}$ where neither the power rule nor the exponential rule applies, iterated exponential towers, and multiplicative functional equations. At the end you should be able to look at a shape and name its rule system before writing anything.
\end{orient}

\section{The rules are typed by shape}

Every elementary expression is a tree. At its root sits one operation: an addition, a multiplication, a division, or a composition. The differentiation rules are typed by that root operation, and the discipline is to identify the root before doing anything else. $\sin(x)\cos(x)$ has multiplication at the root; $\sin(x\cos x)$ has composition at the root and multiplication one level down. The two expressions share every symbol and need entirely different opening moves.

This sounds obvious, and it is the source of most wasted work in the subject. What actually happens is that a solver sees a fraction bar, reaches for the quotient rule reflexively, and only later discovers that the denominator was a single monomial and the whole thing was a sum of three powers. Or a solver sees four factors, applies the two-factor product rule three times in a nested way, and drops a term. Or a solver sees $x^{x}$, recognises neither a power nor an exponential, and produces $x\cdot x^{x-1}=x^{x}$, which is wrong by a factor that happens to be $1+\ln x$.

\begin{keynote}[The four shapes and their systems]
\emph{Nesting}, $f(g(x))$: chain rule, peeled outside in, each derivative evaluated at the remaining inner stack. \emph{Juxtaposition}, $u_{1}u_{2}\cdots u_{n}$: the generalised product rule, one term per factor. \emph{Division}, $u/v$: the quotient rule only if the division is irreducible; otherwise split termwise or write $uv^{-1}$. \emph{Variable exponent}, $f(x)^{g(x)}$: neither the power rule nor the exponential rule, but logarithmic differentiation, which is the only system that covers all four shapes at once. When two or more shapes are present, work from the root down, and re-read the tree after every step.
\end{keynote}

Logarithmic differentiation deserves its position at the end of that list. It is not a fourth rule alongside the others; it is a change of representation that turns multiplication into addition, division into subtraction, and exponentiation into multiplication, so that the product, quotient and power rules all become the sum rule. That is why it is the general solvent, and why the second half of this chapter is devoted to it.

\begin{center}
\begin{tikzpicture}[node distance=6mm and 8mm]
\node[dtest] (a) {Is the exponent variable,\\ or is this a tower?};
\node[dleaf, right=14mm of a] (b) {Logarithmic differentiation:\\ $\ln y=g\ln f$, then $y'=y\cdot(\ln y)'$};
\node[dtest, below=9mm of a] (c) {Four or more nonzero factors,\\ multiplied and divided?};
\node[dleaf, right=14mm of c] (d) {Logarithmic differentiation:\\ $\dfrac{f'}{f}=\sum_{j}\dfrac{u_{j}'}{u_{j}}-\sum_{k}\dfrac{v_{k}'}{v_{k}}$};
\node[dtest, below=9mm of c] (e) {A quotient with one term below,\\ or a sum above?};
\node[dleaf, right=14mm of e] (f) {Split termwise, or write $uv^{-1}$\\ and use product $+$ chain};
\node[dnode, below=9mm of e] (g) {Irreducible two-block quotient: quotient rule.\\ Two or three factors: product rule.\\ Pure nesting: chain rule, outside in.};
\draw[dedge] (a) -- node[dlab,above]{yes} (b);
\draw[dedge] (a) -- node[dlab,left]{no} (c);
\draw[dedge] (c) -- node[dlab,above]{yes} (d);
\draw[dedge] (c) -- node[dlab,left]{no} (e);
\draw[dedge] (e) -- node[dlab,above]{yes} (f);
\draw[dedge] (e) -- node[dlab,left]{no} (g);
\end{tikzpicture}
\end{center}

Two of the four leaves say ``logarithmic differentiation'', which is a fair reflection of how often it is the right answer once an expression is genuinely multiplicative. The bottom box is the ordinary case and it is not a failure: most derivatives really are one product rule or one chain rule, executed cleanly. What the tree prevents is the two expensive misreadings, namely using the quotient rule on a reducible quotient and using a power or exponential rule on a variable-exponent form.

\section{The three rules, used deliberately}

Start with composition, because it is the one rule that never has an alternative. Products can be log-differentiated and quotients can be rewritten, but a genuine nesting must be peeled, and the only question is whether you peel it in a disciplined order.

\tech{Outside-in chain-rule discipline}{easy}
\cue Three or more genuine layers with no available simplification: $\sin\bigl(\cos(1+x^{2})\bigr)$, $\eu^{\tan(\ln x)}$, $\operatorname{arsinh}\sqrt{1+\eu^{3x}}$.
\method Peel from the outermost function inward, never from the inside out:
\[
\dvop{x}f_{1}\bigl(f_{2}(\cdots f_{k}(x))\bigr)=\prod_{j=1}^{k}f_{j}'\bigl(f_{j+1}(\cdots f_{k}(x))\bigr),
\]
where each factor is evaluated at the \emph{remaining inner stack}, not at $x$. Write the layers as a numbered list before differentiating if there are more than three, and substitute the numerical evaluation point only at the very end.

\begin{worked}
Three layers, chosen so the point is kind. $f(x)=\eu^{\tan(\ln x)}$ has layers $\exp$, $\tan$, $\ln$, so
\[
f'(x)=\eu^{\tan(\ln x)}\cdot\sec^{2}(\ln x)\cdot\frac1x .
\]
At $x=1$: $\ln1=0$, $\tan0=0$, $\eu^{0}=1$, $\sec^{2}0=1$, so $f'(1)=1\cdot1\cdot1=1$. Verified: $1.0000000$.

Four layers, with a radical in the middle. $f(x)=\operatorname{arsinh}\sqrt{1+\eu^{3x}}$. With $u=\sqrt{1+\eu^{3x}}$,
\[
f'(x)=\frac{u'}{\sqrt{1+u^{2}}},\qquad u'=\frac{3\eu^{3x}}{2\sqrt{1+\eu^{3x}}},\qquad 1+u^{2}=2+\eu^{3x}.
\]
At $x=0$: $u=\sqrt2$, $u'=\dfrac{3}{2\sqrt2}$, $\sqrt{1+u^{2}}=\sqrt3$, so $f'(0)=\dfrac{3}{2\sqrt6}=\dfrac{\sqrt6}{4}=0.612372$. Verified to seven digits.

And one where the point is deliberately unkind, to show that this is allowed. $f(x)=\sin\bigl(\cos(1+x^{2})\bigr)$ gives $f'(x)=-2x\cos\bigl(\cos(1+x^{2})\bigr)\sin(1+x^{2})$, and at $x=\pi-1$ this is
\[
-2(\pi-1)\cos\Bigl(\cos\bigl(1+(\pi-1)^{2}\bigr)\Bigr)\sin\bigl(1+(\pi-1)^{2}\bigr)=1.979199 .
\]
There is no simplification. The evaluation point contains $\pi$ purely to make you hunt for one.

And, for balance, a nest that should never have been peeled at all. Consider $f(x)=\arctan\bigl(\sinh(\ln x)\bigr)$ for $x>0$. Peeling gives
\[
f'(x)=\frac{1}{1+\sinh^{2}(\ln x)}\cdot\cosh(\ln x)\cdot\frac1x=\frac{\cosh(\ln x)}{x\cosh^{2}(\ln x)}=\frac{1}{x\cosh(\ln x)},
\]
using $1+\sinh^{2}=\cosh^{2}$, and since $\cosh(\ln x)=\tfrac12(x+1/x)$ this is $\dfrac{2}{x^{2}+1}$. So $f'(1)=1$ and $f'(2)=\tfrac25$, both verified to seven digits. The chain rule got there, but the previous chapter's scan would have noticed $\sinh(\ln x)=\tfrac12(x-1/x)$ at the start and reduced the whole thing to a single $\arctan$ of a rational function. Run the scan first; the chain rule is what you use when it comes back empty.
\end{worked}

\begin{trap}
Evaluating an outer derivative at $x$ instead of at the inner block: writing $\cos(x)$ where the formula demands $\cos\bigl(\cos(1+x^{2})\bigr)$. This is the defining error of the chain rule and it survives into every later technique, because Faa di Bruno, implicit differentiation and the inverse-function rule are all chain-rule bookkeeping. The second trap is assuming that an ugly evaluation point \emph{must} trigger a hidden collapse; sometimes the answer is simply transcendental, and the hunt costs more than the arithmetic.
\end{trap}

\begin{seealsob}Composition bookkeeping and functional inversion; the generalised product rule for $n$ factors; logarithmic differentiation of many-factor products.\end{seealsob}

Composition has one rule and no alternative. Multiplication has two: the product rule, which is exact and always available, and logarithmic differentiation, which is faster once the factor count rises but fails where a factor vanishes. It is worth stating the product rule in its $n$-factor form first, because the two-factor form applied repeatedly is where terms get lost.

\tech[Generalised product rule for n factors]{Generalised product rule for $n$ factors}{easy}
\cue Three or more factors multiplied, each nontrivial: $x^{2}\eu^{x}\sin x$, $\ln(x^{2}+3x)\,\eu^{\tan x}\,(x^{3}+1)^{-1}$, $x(x-1)(x-2)(x-3)$.
\method One term per factor, each term differentiating exactly one slot and leaving the rest alone:
\[
\bigl(u_{1}u_{2}\cdots u_{n}\bigr)'=\sum_{j=1}^{n}u_{j}'\prod_{k\neq j}u_{k}.
\]
The result has exactly $n$ terms, which is the checksum: count them before simplifying. When every factor is nonzero at the point, dividing through by $f$ gives $\dfrac{f'}{f}=\sum_{j}\dfrac{u_{j}'}{u_{j}}$, which is logarithmic differentiation and is usually the faster route from four factors up.

\begin{worked}
$f(x)=x^{2}\eu^{x}\sin x$, three factors, so three terms:
\[
f'(x)=2x\,\eu^{x}\sin x+x^{2}\eu^{x}\sin x+x^{2}\eu^{x}\cos x .
\]
At $x=\pi$ the first two terms carry $\sin\pi=0$ and only the third survives: $f'(\pi)=\pi^{2}\eu^{\pi}\cos\pi=-\pi^{2}\eu^{\pi}=-228.3895$. Verified to seven digits. The choice of $\pi$ is the setter arranging for two of the three terms to die, which is a common design and a reason to keep the terms unexpanded until the point is substituted.

The four-factor case, where the product rule beats logarithmic differentiation outright. Let $f(x)=x(x-1)(x-2)(x-3)$ and find $f'(0)$. The generalised rule gives four terms, and three of them retain the factor $x$, which vanishes at $0$. Only the term that differentiates the first slot survives:
\[
f'(0)=1\cdot(0-1)(0-2)(0-3)=-6 .
\]
Verified: $-6.0000000$. Logarithmic differentiation cannot be used here, because $f(0)=0$ and $f'/f$ is undefined at the very point requested. Note how little work the product rule did: at a zero of one factor, $n-1$ of the $n$ terms disappear.

A three-factor case with a chain rule inside one slot, to show the two rules composing. With $u_{1}=\ln(x^{2}+3x)$, $u_{2}=\eu^{\tan x}$, $u_{3}=(x^{3}+1)^{-1}$,
\[
\frac{f'}{f}=\frac{2x+3}{(x^{2}+3x)\ln(x^{2}+3x)}+\sec^{2}x-\frac{3x^{2}}{x^{3}+1},
\]
and at $x=1$ the three contributions are $0.901686$, $3.425519$ and $-1.5$, with $f(1)=3.290024$, giving $f'(1)=9.301566$. Verified by complex-step differentiation.
\end{worked}

\begin{trap}
Nesting the two-factor rule and losing a term. Applying $(uv)'=u'v+uv'$ to $u=u_{1}$, $v=u_{2}u_{3}$ is correct but requires a second application inside, and the standard failure is to differentiate $u_{2}u_{3}$ as $u_{2}'u_{3}'$. Count terms: an $n$-factor product must produce exactly $n$ terms before simplification. The other error is bracketing so that a chain-rule factor generated in one slot silently multiplies a neighbouring term.
\end{trap}

\begin{seealsob}Logarithmic differentiation of many-factor products; when the quotient rule is the wrong tool; outside-in chain-rule discipline.\end{seealsob}

That leaves division, which is the shape most often mishandled, not because the quotient rule is hard but because it is applied when it should not be. The quotient rule is the correct tool for exactly one situation: an irreducible ratio of two genuinely composite blocks. Everything else is cheaper another way, and the difference is not stylistic.

\tech{When the quotient rule is the wrong tool}{easy}
\cue A quotient whose denominator is a monomial, a constant, a single exponential, or a power of one factor; or a quotient whose numerator is a sum; or a quotient of two messy nonzero products.
\method Three cheaper routes, in order of preference. First, split: $\dfrac{a_{1}+\cdots+a_{n}}{g}=\sum_{j}\dfrac{a_{j}}{g}$ whenever $g$ is a single term, then differentiate a sum. Second, rewrite $\dfrac uv=uv^{-1}$ and use the product rule with a chain rule on $v^{-1}$, which avoids squaring the denominator. Third, when $u$ and $v$ are both nonzero products, log-differentiate. Reserve
\[
\left(\frac uv\right)'=\frac{u'v-uv'}{v^{2}}
\]
for a genuinely irreducible two-block quotient such as $\dfrac{x^{2}+1}{x^{2}-1}$ or $\dfrac{\ln x}{\sin x}$.

\begin{worked}
The split, on the standard example:
\[
f(x)=\frac{x^{3}+\sqrt{x}-1}{x^{2}}=x+x^{-3/2}-x^{-2},\qquad f'(x)=1-\tfrac32x^{-5/2}+2x^{-3},
\]
so $f'(4)=1-\tfrac32\cdot\tfrac{1}{32}+\tfrac{2}{64}=1-\tfrac{3}{64}+\tfrac{2}{64}=\dfrac{63}{64}$. Verified: $0.984375$. The quotient rule reaches the same place through a numerator of degree $4$ over $x^{4}$ and a cancellation of $x^{2}$ that must be done by hand, with the half-integer power differentiated inside a bracket.

The negative-power rewrite, where the split is unavailable because the denominator is transcendental. Take $f(x)=\dfrac{\eu^{x}+x^{2}}{\eu^{x}}$. Here the denominator \emph{is} a single term, so the split does apply and is best: $f(x)=1+x^{2}\eu^{-x}$, and
\[
f'(x)=2x\eu^{-x}-x^{2}\eu^{-x}=x(2-x)\eu^{-x},
\]
whence $f'(2)=0$. Verified: $0.0000000$. The quotient rule gives $\dfrac{(\eu^{x}+2x)\eu^{x}-(\eu^{x}+x^{2})\eu^{x}}{\eu^{2x}}$, which is the same number after cancelling $\eu^{x}$ from every term, but the cancellation is invisible until it is performed.

And a case where the quotient rule is genuinely right, so that the technique is a selection rule and not a prohibition. For $f(x)=\dfrac{x^{2}+1}{x^{2}-1}$ neither block factors usefully and neither is a single term, so
\[
f'(x)=\frac{2x(x^{2}-1)-(x^{2}+1)2x}{(x^{2}-1)^{2}}=\frac{-4x}{(x^{2}-1)^{2}},
\]
and $f'(2)=-\dfrac{8}{9}$. Verified: $-0.888889$.
\end{worked}

\begin{trap}
Sign order in the numerator. It is $u'v-uv'$, never $uv'-u'v$, and the mnemonic that survives pressure is that the numerator's derivative comes first, exactly as in the product rule. The second standard error is forgetting to square the denominator. The third is the algebraic one from the previous chapter, splitting over a sum \emph{downstairs}: $\dfrac{1}{a+b}\neq\dfrac1a+\dfrac1b$.
\end{trap}

\begin{seealsob}Algebraic normalisation: roots to powers, quotients to sums; the generalised product rule for $n$ factors; variable-base logarithms.\end{seealsob}

\subsection{One derivative, three rule systems}

It is worth doing a single derivative by every route on offer, because the comparison is this chapter's argument in miniature. Take
\[
f(x)=\frac{x^{2}\sqrt{x+1}}{(x-2)^{3}},\qquad\text{find } f'(3).
\]

\emph{By the quotient rule.} The numerator is itself a product, so a product rule is needed inside: $u=x^{2}\sqrt{x+1}$ gives $u'=2x\sqrt{x+1}+\dfrac{x^{2}}{2\sqrt{x+1}}$, and $v=(x-2)^{3}$ gives $v'=3(x-2)^{2}$. Then $f'=\dfrac{u'v-uv'}{v^{2}}$ with a sixth-degree denominator that must be cancelled down to a cube. At $x=3$: $u=9\cdot2=18$, $u'=6\cdot2+\tfrac94=\tfrac{57}{4}$, $v=1$, $v'=3$, so $f'(3)=\tfrac{57}{4}-54=-\tfrac{159}{4}$. Correct, and four places to lose a factor.

\emph{By the product rule with a negative power.} Write $f=x^{2}(x+1)^{1/2}(x-2)^{-3}$, a three-factor product, and use the generalised rule:
\[
f'=2x(x+1)^{1/2}(x-2)^{-3}+\tfrac12x^{2}(x+1)^{-1/2}(x-2)^{-3}-3x^{2}(x+1)^{1/2}(x-2)^{-4}.
\]
At $x=3$ the three terms are $12$, $\tfrac94$ and $-54$, summing to $-\tfrac{159}{4}$. No denominator was ever squared.

\emph{By logarithmic differentiation.} Three nonzero factors and a specific point, so
\[
\frac{f'}{f}=\frac{2}{x}+\frac{1}{2(x+1)}-\frac{3}{x-2},
\]
which at $x=3$ is $\tfrac23+\tfrac18-3=-\tfrac{53}{24}$. With $f(3)=\dfrac{9\cdot2}{1}=18$,
\[
f'(3)=18\cdot\left(-\frac{53}{24}\right)=-\frac{159}{4}=-39.75 .
\]
Verified by complex-step differentiation: $-39.750000$. Three routes, one answer, and a clear ranking: the third is two lines and involves no differentiation of anything but $\ln$ of linear and radical factors; the second is three lines and needs no cancellation; the first is correct and costs a page.

The ranking flips if the question changes. Asked for $f'$ as a \emph{function} in simplified form, logarithmic differentiation ends with a multiplication by $f$ that reassembles the whole expression, and the product-rule route with a common factor pulled out is cleaner. Asked for $f'(2)$, where $f$ is undefined, none of the three applies and the question is about a pole. Read what is asked before choosing.

One more piece of rule discipline before the solvent. The chain rule is usually read forwards, from the composition to its derivative. Read backwards it answers a different and commonly set question: you are told about a composition and asked about one of its ingredients.

\tech{Composition bookkeeping and functional inversion}{medium}
\cue You are given $q$ and $p\circ q$ but asked for $p'$ at an \emph{output} value; or you are asked for $(f\circ g^{-1})'(c)$ with $g^{-1}$ not available in closed form.
\method The chain rule solved for the unknown factor. From $(p\circ q)'(a)=p'\bigl(q(a)\bigr)q'(a)$: first solve $q(a)=\text{target}$ for $a$, then
\[
p'(\text{target})=\frac{(p\circ q)'(a)}{q'(a)}.
\]
For the inverse form, find $a$ with $g(a)=c$, which is a root-finding step and not a differentiation step, and then $(f\circ g^{-1})'(c)=\dfrac{f'(a)}{g'(a)}$. The whole difficulty is bookkeeping: knowing which argument each derivative is evaluated at.

\begin{worked}
Given $q(x)=2x+1$ and $(p\circ q)(x)=4x^{2}-2x+3$, find $p'(12)$. The target is an output of $q$, so solve $q(a)=12$: $2a+1=12$ gives $a=\tfrac{11}{2}$. Then $(p\circ q)'(x)=8x-2$, so $(p\circ q)'\bigl(\tfrac{11}{2}\bigr)=44-2=42$, and $q'\equiv2$. Hence
\[
p'(12)=\frac{42}{2}=21 .
\]
As a check, $p$ is available here: $p(u)=(p\circ q)\bigl(\tfrac{u-1}{2}\bigr)=(u-1)^{2}-(u-1)+3=u^{2}-3u+5$, so $p'(u)=2u-3$ and $p'(12)=21$.

Now the inverse form, where $g^{-1}$ is not worth computing. Let $f(x)=\ln(x^{2}+2x+4)$ and $g(x)=\dfrac{4}{1+\eu^{-2x}}$, and find $(f\circ g^{-1})'(2)$. First locate $a$ with $g(a)=2$: the logistic is $2$ at its midpoint, where $\eu^{-2a}=1$, so $a=0$. Then
\[
f'(x)=\frac{2x+2}{x^{2}+2x+4}\ \Longrightarrow\ f'(0)=\frac24=\frac12,\qquad
g'(x)=\frac{8\eu^{-2x}}{(1+\eu^{-2x})^{2}}\ \Longrightarrow\ g'(0)=\frac84=2,
\]
so $(f\circ g^{-1})'(2)=\dfrac{1/2}{2}=\dfrac14$. Verified by complex-step differentiation of both functions at $0$: $f'(0)=0.5000000$ and $g'(0)=2.0000000$.
\end{worked}

\begin{trap}
Evaluating $p'$ at $a$ rather than at $q(a)$. In the first example that would give $p'\bigl(\tfrac{11}{2}\bigr)=8$ instead of $p'(12)=21$, and nothing in the arithmetic complains. The companion error is dividing by $q'$ evaluated at the target instead of at $a$; here $q'$ is constant so the error hides, which is exactly why the example is a poor test of understanding and a good place to build the habit. Always write the argument of every derivative explicitly.
\end{trap}

\begin{seealsob}Outside-in chain-rule discipline; the general power rule for $f(x)^{g(x)}$; variable-base logarithms.\end{seealsob}

\section{Logarithmic differentiation as the general solvent}

The three rules and their selection guidance are now complete, and between them they differentiate any expression built from sums, products, quotients and compositions. What they do not touch is the fourth shape, the variable exponent, and they handle the third shape, division by a many-factor product, only expensively. Both gaps are closed by one manoeuvre, which is the subject of the rest of the chapter.

Here is the manoeuvre in three lines. Given $y=f(x)$ with $f$ nonzero near the point, take logarithms of both sides, differentiate implicitly using $\dvop{x}\ln y=\dfrac{y'}{y}$, and multiply back:
\[
\ln y=\ln f(x)\quad\Longrightarrow\quad \frac{y'}{y}=\bigl(\ln f\bigr)'(x)\quad\Longrightarrow\quad y'=f(x)\cdot\bigl(\ln f\bigr)'(x).
\]
Nothing here is deep. What makes it a solvent is what the logarithm does to structure on the way in: products become sums, quotients become differences, powers become products, and towers lose a storey. Every one of those is a reduction in the complexity the differentiation rules have to face. The quantity $f'/f$ has a name, the logarithmic derivative, and it is worth thinking of as an operator $\mathcal{L}[f]=f'/f$ with the property $\mathcal{L}[fg]=\mathcal{L}[f]+\mathcal{L}[g]$: it converts the multiplicative structure of functions into the additive structure of their derivatives.

\begin{keynote}[$\mathcal{L}$ turns multiplication into addition]
The logarithmic derivative $\mathcal{L}[f]=f'/f$ satisfies
\[
\mathcal{L}[fg]=\mathcal{L}[f]+\mathcal{L}[g],\quad
\mathcal{L}[f/g]=\mathcal{L}[f]-\mathcal{L}[g],\quad
\mathcal{L}[f^{a}]=a\,\mathcal{L}[f],\quad
\mathcal{L}[c]=0,\quad
\mathcal{L}[cf]=\mathcal{L}[f].
\]
Those five lines are the whole method. The first says a product rule becomes a sum; the second says a quotient rule becomes a difference; the third says a power rule becomes a multiplication; the fourth and fifth say constants are invisible, which is why $\mathcal{L}$ never carries the clutter that $f'$ does. Note also $\mathcal{L}[f]=\bigl(\ln\abs{f}\bigr)'$, so the operator is well defined wherever $f\neq0$, sign notwithstanding.
\end{keynote}

Two conditions govern its use. First, $f$ must be nonzero at and near the point, since $f'/f$ is otherwise undefined; at a zero, fall back on the product rule, which handles zeros without complaint. Second, when factors are negative you are really working with $\ln\abs{f}$, whose derivative is still $f'/f$, so the final formula is unaffected, but any intermediate reasoning about $\ln$ of a negative factor is not licensed.

\tech{Logarithmic differentiation of many-factor products and quotients}{core}
\cue Four or more nonzero factors multiplied and divided, especially many linear factors, with a \emph{specific} evaluation point requested.
\method Apply $\mathcal{L}$ to the whole expression:
\[
\frac{f'}{f}=\sum_{j}\frac{u_{j}'}{u_{j}}-\sum_{k}\frac{v_{k}'}{v_{k}} .
\]
Evaluate that sum of reciprocals at the point, then multiply by $f$ at the point. For a product of linear factors it collapses to $\dfrac{f'}{f}=\sum_{j}\dfrac{1}{x-r_{j}}$, a sum of simple reciprocals with no differentiation left in it at all.

\begin{worked}
\[
f(x)=\frac{(x-1)(x-2)(x-3)}{(x-4)(x-5)},\qquad \text{find } f'(0).
\]
All five factors are nonzero at $0$, so
\[
\frac{f'(0)}{f(0)}=\frac{1}{-1}+\frac{1}{-2}+\frac{1}{-3}-\frac{1}{-4}-\frac{1}{-5}=-\frac{11}{6}+\frac{9}{20}=-\frac{83}{60},
\]
and $f(0)=\dfrac{(-1)(-2)(-3)}{(-4)(-5)}=-\dfrac{3}{10}$. Hence
\[
f'(0)=\left(-\frac{3}{10}\right)\left(-\frac{83}{60}\right)=\frac{83}{200}=0.415 .
\]
Verified by complex-step differentiation: $0.4150000$. The quotient rule on this expression requires a product rule on a cubic numerator, a product rule on a quadratic denominator, and a fourth-degree denominator to square.

Scale it up and the advantage becomes absurd. For $f(x)=\prod_{k=1}^{10}(x-k)$ at $x=0$,
\[
\frac{f'(0)}{f(0)}=-\sum_{k=1}^{10}\frac1k=-H_{10}=-\frac{7381}{2520},\qquad f(0)=(-1)^{10}\,10!=3\,628\,800,
\]
so $f'(0)=-3\,628\,800\cdot\dfrac{7381}{2520}=-1440\cdot7381=-10\,628\,640$. Verified to nine significant figures. The product rule would produce ten terms of nine factors each.
\end{worked}

\begin{trap}
Using it where a factor vanishes at the evaluation point. If $f(x)=x(x-1)(x-2)(x-3)$ and you want $f'(0)$, then $f(0)=0$ and $f'/f$ is undefined, so the method reports nothing while the answer is a perfectly ordinary $-6$. The rule is: a zero of $f$ at the point sends you back to the generalised product rule, where $n-1$ of the $n$ terms vanish and the computation is a one-liner. The second, smaller trap is worrying about $\ln$ of a negative factor. Work with $\ln\abs{u_{j}}$ throughout; the formula $f'/f=\sum u_{j}'/u_{j}$ is unchanged.
\end{trap}

\begin{seealsob}The generalised product rule for $n$ factors; the logarithmic derivative as an operator on functional equations; the general power rule for $f(x)^{g(x)}$.\end{seealsob}

Products were the easy case, since the product rule was always available as an alternative. The variable-exponent form has no alternative at all, and it is the place where students most reliably produce a confidently wrong answer, because two familiar rules both almost apply.

\tech[General power rule for f(x) to the g(x)]{The general power rule: $f(x)^{g(x)}$}{medium}
\cue A variable base raised to a variable exponent: $(\sin x)^{\tan x}$, $x^{\sin x}$, $(1+x)^{\ln x}$, $(3\sin x)^{1/(x^{2}+1)}$, $\Gamma(x)^{W(x)}$.
\method Neither elementary rule applies, and the reason is worth stating precisely. The power rule $\dvop{x}u^{n}=nu^{n-1}u'$ is proved for constant $n$ and its proof uses the binomial theorem with a fixed exponent. The exponential rule $\dvop{x}a^{v}=a^{v}\ln a\cdot v'$ is proved for constant base $a$ and its proof uses $a^{v}=\eu^{v\ln a}$ with $\ln a$ a number. When both vary, write $y=\eu^{g\ln f}$ or equivalently take logarithms, $\ln y=g\ln f$, and differentiate:
\[
\frac{y'}{y}=g'\ln f+\frac{gf'}{f},\qquad\text{so}\qquad y'=f^{g}\left(g'\ln f+\frac{gf'}{f}\right).
\]
The two terms are exactly the two wrong answers: the first is what you get pretending the base is constant, the second is what you get pretending the exponent is constant. Both are always present.

\begin{keynote}[The two wrong answers to $x^{x}$]
Ask three solvers for $\dvop{x}x^{x}$ and you will get $x\cdot x^{x-1}=x^{x}$ from the one who saw a power, $x^{x}\ln x$ from the one who saw an exponential, and $x^{x}(\ln x+1)$ from the one who saw neither. The third is right, and it is exactly the sum of the first two. That is not a coincidence: $g'\ln f$ is the exponential contribution and $gf'/f$ is the power contribution, and the general formula says to add them. Whenever you catch yourself about to apply one familiar rule to a variable-exponent form, write both terms instead.
\end{keynote}

\begin{worked}
$y=(\sin x)^{\tan x}$ at $x=\pi/4$. Here $f=\sin x$, $g=\tan x$, so
\[
\frac{y'}{y}=\sec^{2}x\,\ln\sin x+\tan x\cdot\frac{\cos x}{\sin x}=\sec^{2}x\,\ln\sin x+1,
\]
the second term collapsing because $\tan x\cot x=1$. At $x=\pi/4$: $\sec^{2}=2$, $\ln\sin(\pi/4)=\ln\tfrac{1}{\sqrt2}=-\tfrac12\ln2$, and $y=(\tfrac{1}{\sqrt2})^{1}=\tfrac{1}{\sqrt2}$. So
\[
y'\left(\frac\pi4\right)=\frac{1}{\sqrt2}\bigl(1-\ln 2\bigr)=0.216978 .
\]
Verified to seven digits.

A cleaner design, where the point kills one term outright. For $y=x^{\sin x}$ at $x=\pi$: $\dfrac{y'}{y}=\cos x\,\ln x+\dfrac{\sin x}{x}$, and at $\pi$ the second term vanishes while $y(\pi)=\pi^{0}=1$. Hence $y'(\pi)=-\ln\pi=-1.144730$. Verified. Setters engineer this constantly: choose the point so that $\ln f=0$ (that is, $f=1$) or $g=0$, and one of the two terms disappears.

The same design with a special function. For $f(x)=\Gamma(x)^{W(x)}$ at $x=1$: $f(1)=\Gamma(1)^{W(1)}=1$, and
\[
\frac{f'(1)}{f(1)}=W'(1)\ln\Gamma(1)+W(1)\,\frac{\Gamma'(1)}{\Gamma(1)}=W'(1)\cdot0+\Omega\cdot\psi(1)=-\gamma\,\Omega,
\]
using $\Gamma(1)=1$, $\psi(1)=-\gamma$, and $W(1)=\Omega=0.5671433$, the omega constant. So $f'(1)=-\gamma\Omega=-0.3273640$. The unknown quantity $W'(1)$ never had to be computed, because $\ln\Gamma(1)=0$ deleted it.

A case where the other term dies instead. For $y=(3\sin x)^{1/(x^{2}+1)}$ at $x=\pi/2$, the exponent is $g=(x^{2}+1)^{-1}$ and the base is $f=3\sin x$, so
\[
\frac{y'}{y}=\frac{-2x}{(x^{2}+1)^{2}}\ln(3\sin x)+\frac{1}{x^{2}+1}\cdot\cot x .
\]
At $x=\pi/2$ the cotangent vanishes, $\sin x=1$, and $y=3^{1/(\pi^{2}/4+1)}=1.372783$, giving
\[
y'\!\left(\frac\pi2\right)=1.372783\cdot\frac{-\pi\ln 3}{\bigl(\pi^{2}/4+1\bigr)^{2}}=-0.394083 .
\]
Verified to seven digits. Here it was the base's contribution that died and the exponent's that survived, the mirror image of the $\Gamma^{W}$ example, and in both cases the point was chosen for exactly that purpose.
\end{worked}

\begin{trap}
Writing $g f^{g-1}f'$, which is valid only for constant $g$, or $f^{g}\ln f\cdot g'$, which is valid only for constant $f$. Test both against $y=x^{x}$, whose true derivative is $x^{x}(\ln x+1)$: the first gives $x\cdot x^{x-1}=x^{x}$, the second gives $x^{x}\ln x$, and their sum is the correct answer, which is not a coincidence but the two terms of the formula. A subtler trap: $\ln f$ requires $f>0$, so $(-2)^{x}$ and $(\cos x)^{x}$ near $\cos x<0$ need care, and the expression may not be real at all.
\end{trap}

\begin{seealsob}Power towers and iterated exponentials; exponent arithmetic before log-differentiating; logarithmic differentiation of many-factor products.\end{seealsob}

One variable exponent needs one pass of logarithmic differentiation. Two stacked variable exponents need two, and the only new difficulty is bookkeeping: knowing what counts as the base, what counts as the exponent, and which way the tower associates.

\tech{Power towers and iterated exponentials}{hard}
\cue Two or more stacked variable exponents: $x^{x^{x}}$, $x^{x^{\sin x}}$, $(x^{x})^{\ln x}$, or a finite tower of any height.
\method Log-differentiate \emph{outermost first}, treating the entire inner tower as a single block $g$. For $f=x^{g}$ we get $\ln f=g\ln x$ and
\[
\frac{f'}{f}=g'\ln x+\frac gx ,
\]
which is the general power rule with $f=x$. Now $g$ is itself a tower and needs its own pass. Recurse until the exponent is elementary. Before any of this, settle the association: towers are read top-down, $x^{x^{x}}=x^{(x^{x})}$, and in particular $(x^{x})^{\ln x}=x^{x\ln x}$ is a different function from $x^{(x^{\ln x})}$.

\begin{worked}
The bracketed case first, because it collapses. $f(x)=(x^{x})^{\ln x}=x^{x\ln x}=\eu^{x(\ln x)^{2}}$, so the exponent is elementary and one pass suffices:
\[
\frac{f'}{f}=(\ln x)^{2}+x\cdot2\ln x\cdot\frac1x=(\ln x)^{2}+2\ln x .
\]
At $x=\eu$: $(\ln \eu)^{2}+2\ln \eu=1+2=3$ and $f(\eu)=\eu^{\eu}$, so $f'(\eu)=3\eu^{\eu}=45.46279$. Verified to seven digits.

Now the genuine two-storey tower, $f(x)=x^{x^{x}}$. Outermost pass with $g=x^{x}$:
\[
\frac{f'}{f}=g'\ln x+\frac gx .
\]
Inner pass: $g=x^{x}$ has $\dfrac{g'}{g}=\ln x+1$, so $g'=x^{x}(\ln x+1)$. Substituting,
\[
f'(x)=x^{x^{x}}\left(x^{x}(\ln x+1)\ln x+x^{x-1}\right).
\]
At $x=1$: $\ln1=0$ kills the first term, $x^{x-1}=1$, and $f(1)=1$, so $f'(1)=1$. At $x=2$: $x^{x}=4$, $\ln2=0.693147$, $f(2)=2^{4}=16$, and
\[
f'(2)=16\bigl(4(1+\ln2)\ln2+2\bigr)=64\ln2\,(1+\ln2)+32=107.11041 .
\]
Both verified to seven digits.
\end{worked}

\begin{trap}
Reading the tower with the wrong association. $x^{x^{2}}=x^{(x^{2})}=\eu^{x^{2}\ln x}$, whereas $(x^{x})^{2}=x^{2x}=\eu^{2x\ln x}$: different functions, different derivatives, identical-looking notation if the brackets are dropped. The second trap is stopping after one pass while $g'$ is still an unevaluated symbol; the answer is not finished until every storey has been differentiated. The third is losing the term $g/x$, which is the contribution of the base and is easy to drop when the exponent looks like all the action.
\end{trap}

\begin{seealsob}The general power rule for $f(x)^{g(x)}$; exponent arithmetic before log-differentiating; self-similar fixed-point resolution.\end{seealsob}

Towers are one way to make an exponent complicated. The other is to make it long: a product of many powers of $x$, where the exponents must be summed before anything else happens. This is a chapter-eleven move applied inside a chapter-twelve technique, and it is worth its own block because the collapse is invisible if you start differentiating.

\tech{Exponent arithmetic before log-differentiating}{medium}
\cue A finite product of powers of the same base, or stacked radicals of such powers: $\left(\prod_{i=1}^{n}x^{i}\right)^{2}$, or a chain like $\dfrac{x^{10x^{2}}}{\sqrt{x^{x^{3}}}}\cdot\dfrac{\sqrt[3]{x^{x^{3}}}}{x^{5x^{3}/2}}$.
\method All the work is in the exponent. Collapse $\prod_{i}x^{e_{i}}=x^{\sum_{i}e_{i}}$ using whatever summation the exponents invite: the arithmetic series $\sum_{i=1}^{n}i=\tfrac{n(n+1)}{2}$, a geometric sum, or a telescoping cancellation. Only then apply the general power rule. A useful special value: at $x=1$ every such $f$ has $f(1)=1$ and $\ln1=0$, so $f'(1)=E(1)$, the exponent evaluated at $1$.

\begin{worked}
Take $f(x)=\left(\prod_{i=1}^{n}x^{i}\right)^{2}$ with the index bound set equal to the variable, as these puzzles are usually posed. The inner product is $x^{1+2+\cdots+n}=x^{n(n+1)/2}$, the outer square doubles the exponent, and with $n=x$,
\[
f(x)=x^{x(x+1)}=x^{x^{2}+x}.
\]
Now one pass of the general power rule with $E(x)=x^{2}+x$:
\[
\frac{f'}{f}=E'(x)\ln x+\frac{E(x)}{x}=(2x+1)\ln x+x+1 .
\]
At $x=2$: $E(2)=6$, so $f(2)=2^{6}=64$, and $f'/f=5\ln2+3$. Hence
\[
f'(2)=64\bigl(3+5\ln2\bigr)=413.8071 .
\]
Verified to seven digits. At $x=1$ the shortcut applies: $f(1)=1$ and $f'(1)=E(1)=2$, confirmed numerically as $2.0000000$.

The stacked-radical chain works the same way. Writing $\sqrt{x^{x^{3}}}=x^{x^{3}/2}$ and $\sqrt[3]{x^{x^{3}}}=x^{x^{3}/3}$, the four displayed factors contribute exponents $10x^{2}$, $-\tfrac12x^{3}$, $+\tfrac13x^{3}$ and $-\tfrac52x^{3}$, so $E(x)=10x^{2}-\tfrac{8}{3}x^{3}$ and one pass finishes it. The point is that no differentiation happened until the exponent arithmetic was complete.
\end{worked}

\begin{trap}
Differentiating factor by factor as a giant product and never noticing that the exponents telescope. With $n$ factors that is $n$ applications of the general power rule and $n$ terms to add, in place of one arithmetic series. The other error is forgetting that an outer square or root multiplies \emph{every} exponent: $\left(x^{n(n+1)/2}\right)^{2}=x^{n(n+1)}$, not $x^{n(n+1)/2+2}$ and not $x^{2n(n+1)/2\cdot 2}$.
\end{trap}

\begin{seealsob}Power towers and iterated exponentials; the general power rule for $f(x)^{g(x)}$; sum the hidden series, product, or continued fraction first.\end{seealsob}

Logarithms have appeared so far as a tool applied to the outside of an expression. They can also appear inside it, and when the \emph{base} of a logarithm depends on $x$ the change-of-base formula is compulsory rather than optional.

\tech[Variable-base logarithms]{Variable-base logarithms $\log_{u(x)}v(x)$}{medium}
\cue A logarithm whose base is not a constant: $\log_{x}2$, $\log_{x}(x^{2}+1)$, $\log_{x^{2}+3x}(x^{3}+\sin x)$.
\method Change base first, to natural logarithms, and only then differentiate:
\[
\log_{u}v=\frac{\ln v}{\ln u}\quad\Longrightarrow\quad
\dvop{x}\log_{u}v=\frac{\dfrac{v'}{v}\ln u-\dfrac{u'}{u}\ln v}{(\ln u)^{2}} .
\]
This is one of the few places where the quotient rule is genuinely the right tool: after the change of base the expression is an irreducible ratio of two transcendental blocks, and neither splitting nor a negative power helps.

\begin{worked}
A short one to fix the shape. $f(x)=\log_{x}2=\dfrac{\ln2}{\ln x}$ has a constant numerator, so
\[
f'(x)=\ln2\cdot\left(-\frac{1}{(\ln x)^{2}}\right)\cdot\frac1x=-\frac{\ln2}{x(\ln x)^{2}},
\]
and $f'(\eu)=-\dfrac{\ln2}{\eu}=-0.2549946$. Verified.

The full two-term case. Let $f(x)=\log_{x^{2}+3x}\bigl(x^{3}+\sin x\bigr)$ and evaluate $f'(3)$. Here $u=x^{2}+3x$ with $u(3)=18$ and $u'(3)=2x+3=9$; $v=x^{3}+\sin x$ with $v(3)=27+\sin3=27.141120$ and $v'(3)=3x^{2}+\cos x=27+\cos3=26.010007$. Then
\[
\frac{v'}{v}=0.958331,\quad \ln u=\ln18=2.890372,\quad \frac{u'}{u}=\frac12,\quad \ln v=3.301207,
\]
so the numerator is $0.958331\times2.890372-0.5\times3.301207=2.769737-1.650604=1.119133$, and dividing by $(\ln18)^{2}=8.354250$ gives
\[
f'(3)=0.1339905 .
\]
Verified by complex-step differentiation of $\ln(x^{3}+\sin x)/\ln(x^{2}+3x)$: $0.1339905$.
\end{worked}

\begin{trap}
Treating the base as constant and writing only $\dfrac{v'}{v\ln u}$. In the example above that term alone gives $0.331565$, against a true value of $0.133990$: the neglected second term is more than half the answer. The error is easy to make because $\log_{a}$ with a numeric $a$ is so much more familiar, and because the notation puts the base in small type.
\end{trap}

\begin{seealsob}When the quotient rule is the wrong tool; logarithmic differentiation of many-factor products; the general power rule for $f(x)^{g(x)}$.\end{seealsob}

The last technique in the chapter uses the logarithmic derivative not as a computational device but as an operator with an algebraic property, applied to a function that is never written down. This is the most abstract use and the most powerful: it solves problems where no formula for $f$ exists.

\tech{The logarithmic derivative as an operator on functional equations}{hard}
\cue $f$ is given only through a multiplicative functional equation, such as $f(x+1)=x\,f(x)$ or $f(2x)=f(x)^{2}g(x)$, and the question involves $g=f'/f$ rather than $f$ itself.
\method Apply $\mathcal{L}[f]=f'/f$ to both sides. Because $\mathcal{L}[fg]=\mathcal{L}[f]+\mathcal{L}[g]$ and $\mathcal{L}[c]=0$ for constants, a multiplicative recurrence in $f$ becomes an \emph{additive} recurrence in $g$, which telescopes. From $f(x+1)=x f(x)$, taking logarithms gives $\ln f(x+1)=\ln x+\ln f(x)$ and differentiating gives
\[
g(x+1)=g(x)+\frac1x ,
\]
which is the digamma recurrence $\psi(x+1)=\psi(x)+1/x$. Never identify $f$; just telescope.

\begin{worked}
Suppose $f$ satisfies $f(x+1)=x f(x)$ on $x>0$ and $g=f'/f$. Compute $12\bigl(g(5)-g(1)\bigr)$.

Applying $\mathcal{L}$ to both sides gives $g(x+1)-g(x)=1/x$, so telescoping from $1$ to $5$,
\[
g(5)-g(1)=\sum_{k=1}^{4}\frac1k=1+\frac12+\frac13+\frac14=\frac{25}{12},
\]
and therefore $12\bigl(g(5)-g(1)\bigr)=25$. Verified numerically against the digamma function, which is one solution of the recurrence: $\psi(5)-\psi(1)=2.083313$ from a truncated series, against $25/12=2.083333$.

Note what was not needed. No formula for $f$, no value of $f$ anywhere, no convergence argument, and no special-function theory. The multiplicative recurrence was converted into an additive one and the additive one was summed. The same move handles $f(2x)=f(x)^{2}$, which gives $2g(2x)=2g(x)$, that is $g(2x)=g(x)$, so the logarithmic derivative is invariant under doubling.

The general telescope is worth recording, since it is the shape every such problem takes. For any integer $n\geq1$,
\[
g(n)=g(1)+\sum_{k=1}^{n-1}\frac1k=g(1)+H_{n-1},
\]
so the entire integer behaviour of $g$ is fixed by one constant. If $f$ is normalised to be $\Gamma$, that constant is $g(1)=\psi(1)=-\gamma$ and the formula reads $\psi(n)=-\gamma+H_{n-1}$, the standard evaluation of the digamma function at integers. But note the direction of the logic: the recurrence gave every \emph{difference} for free, and the normalisation was needed only to fix the additive constant. Questions that ask for a difference, as this one did, never need it.
\end{worked}

\begin{trap}
Trying to identify $f$ explicitly as $\Gamma$. The functional equation $f(x+1)=xf(x)$ does not determine $f$: multiply any solution by a period-one function such as $\eu^{\sin 2\pi x}$ and you get another. Pinning down $\Gamma$ requires an extra hypothesis, logarithmic convexity in the Bohr and Mollerup theorem, which the problem has not given you and does not need. Work with the difference equation and stop.
\end{trap}

\begin{seealsob}Logarithmic differentiation of many-factor products; the general power rule for $f(x)^{g(x)}$; composition bookkeeping and functional inversion.\end{seealsob}

A closing word on when \emph{not} to use the solvent. Logarithmic differentiation costs a logarithm, a derivative, and a multiplication back by $f$, and that overhead is only repaid when the expression is genuinely multiplicative. On a sum it is useless, since $\ln(u+v)$ has no expansion. On a two-factor product it is slower than the product rule. At a zero of $f$ it is undefined. And when the question asks for $f'$ as a function rather than at a point, the final multiplication by $f$ reassembles all the mess you just dissolved, which is sometimes worse than never dissolving it. The technique earns its name on many-factor products evaluated at a point, on variable exponents where nothing else works, and on functional equations where $f$ is unavailable. Those three cases are most of what this chapter has been about.

Two habits carry all of it. The first is naming the root operation before differentiating, so that the rule is chosen rather than defaulted to. The second is asking, whenever the expression is multiplicative, whether taking a logarithm would make the problem additive; if it would, and if $f$ does not vanish at the point, it almost always should be taken. Between them they eliminate the two failure modes this chapter opened with: algebra that never ends, and a confident answer that is missing one of its two terms.

\section{Recognition matrix}
\begin{recog}{@{}p{0.30\textwidth}p{0.36\textwidth}p{0.28\textwidth}@{}}
\rowcolor{mist}\mh{If you see} & \mh{Reach for} & \mh{If that fails} \\
\midrule
\endhead
Three or more genuine layers, no simplification & Chain rule outside in; each factor at the inner stack & List the layers before differentiating \\
An ugly evaluation point in a nested composition & Peel first, substitute last & Accept a transcendental answer \\
Three or more nontrivial factors & Generalised product rule: exactly $n$ terms & Check the term count as a checksum \\
A product evaluated where one factor vanishes & Product rule: $n-1$ terms die & Never $f'/f$; it is undefined there \\
A quotient with a monomial denominator & Split termwise into powers, then the power rule & Rewrite as $uv^{-1}$ \\
A quotient with a sum on top & Split; you may split above, never below & Product $+$ chain on $uv^{-1}$ \\
An irreducible two-block quotient & Quotient rule $\dfrac{u'v-uv'}{v^{2}}$; square the denominator & Log-differentiate if both blocks are products \\
Given $p\circ q$ and $q$, asked for $p'$ at an output & Solve $q(a)=$ target, then divide by $q'(a)$ & Recover $p$ explicitly and check \\
$(f\circ g^{-1})'(c)$ & Find $a$ with $g(a)=c$; answer is $f'(a)/g'(a)$ & Inverse-function rule on $g$ alone \\
Four or more nonzero factors, a point requested & $\dfrac{f'}{f}=\sum\dfrac{u_j'}{u_j}-\sum\dfrac{v_k'}{v_k}$, then multiply by $f$ & Product rule if any factor vanishes \\
Many linear factors & $\dfrac{f'}{f}=\sum_{j}\dfrac{1}{x-r_{j}}$; no differentiation left & Harmonic sums appear at integer points \\
$f(x)^{g(x)}$ with both variable & $y'=f^{g}\bigl(g'\ln f+gf'/f\bigr)$; both terms always & Pick the point that kills one term \\
A stacked tower $x^{x^{x}}$ & Log-differentiate outermost first, then recurse & Fix the association before starting \\
A long product of powers of $x$ & Sum the exponents first, then one power-rule pass & At $x=1$, $f'(1)=E(1)$ \\
$\log_{u(x)}v(x)$ & Change base to $\ln v/\ln u$, then the quotient rule & Both terms matter; the base is not inert \\
$f(x+1)=x f(x)$, or any multiplicative recurrence & Apply $\mathcal{L}[f]=f'/f$; it becomes additive & Telescope; do not identify $f$ \\
\end{recog}
